\section{实验设置：完整细节}
\label{app:setup}

本附录在 \S\ref{sec:experiments:setup} 的精简版设置基础上加以扩展，给出形式化的指标定义与运行时基础设施。

\paragraph{种子智能体。}
种子配置记为 \textbf{NexAU\textsubscript{0}}，它是构建在 \S\ref{sec:method:harness} 所述 NexAU 框架之上的一个简单 code agent，只向模型暴露 \texttt{bash} 工具，没有技能、没有中间件、也没有长期记忆。AHE 外层循环的每一轮迭代都编辑这个工作区，因此所报告的全部增益都是以 NexAU\textsubscript{0} 作为共同起点来度量的。

\paragraph{运行时基础设施。}
所有运行都使用 \S\ref{sec:method:harness} 的 NexAU 框架来实例化 coding agent。Harbor 负责派发任务、隔离每一次 rollout，并验证通过/失败，单任务超时为 3600 秒。每一次 rollout 都在一个全新的 E2B 远程沙箱中运行，因此 shell 的副作用不会在任务之间泄漏。\texttt{InMemoryTracer} 记录轨迹并将其镜像到 Langfuse。Agent Debugger 以并发度 16 执行，单任务超时为 600 秒。

\paragraph{超参数。}
表~\ref{tab:hparams} 汇总了四个智能体与外层循环的工作点，取自参考运行的配置快照。

\begin{table}[h]
    \centering
    \small
    \setlength{\tabcolsep}{6pt}
    \renewcommand{\arraystretch}{1.05}
    \caption{Terminal-Bench~2 上参考 AHE 运行的超参数，取自产生 \S\ref{sec:experiments:main} 中 AHE 数值的那次运行的配置快照。四个智能体都以 GPT-5.4 作为底层模型，仅在推理档位、采样以及各组件的限额上有所不同。}
    \label{tab:hparams}
    \begin{tabular}{lll}
        \toprule
        模块 & 超参数 & 取值 \\
        \midrule
        Code Agent
            & 底层模型              & GPT-5.4 \\
            & 推理强度           & high \\
            & 温度                & 0.7 \\
            & Top-$p$                    & 0.95 \\
            & 最大上下文                & 200{,}000 tokens \\
            & 单轮最大生成长度    & 32{,}000 tokens \\
            & 最大模型轮数            & 300 \\
        \midrule
        Agent Debugger
            & 底层模型              & GPT-5.4 \\
            & 并发任务数           & 16 \\
            & 单任务超时           & 600\,s \\
            & 重试次数            & 3 \\
            & 重试退避             & 2\,s \\
        \midrule
        Evolve Agent
            & 底层模型              & GPT-5.4 \\
            & 推理强度           & xhigh \\
            & 温度                & 0.7 \\
            & 最大上下文                & 200{,}000 tokens \\
            & 单轮最大生成长度    & 32{,}000 tokens \\
            & 最大模型轮数            & 500 \\
            & 技能包             & 3 \\
            & 上下文压缩阈值 & 0.75 \\
        \midrule
        Explore Agent
            & 底层模型              & GPT-5.4 \\
            & 推理强度           & xhigh \\
            & 墙钟时间预算          & 60\,min \\
            & 网络来源数                & 10 \\
            & 代码来源数               & 1 \\
        \midrule
        外层循环 / 调度器
            & 迭代轮数 $T$             & 10 \\
            & 每任务 rollout 数 $k$      & 2 \\
            & 并发 rollout 数        & 96 \\
            & 沙箱                    & E2B remote \\
            & 沙箱生命周期            & 3{,}600\,s \\
            & 数据集                     & Terminal-Bench~2, 89 个任务 \\
        \bottomrule
    \end{tabular}
\end{table}


\paragraph{Terminal-bench 难度标签。}
官方 terminal-bench-2 排行榜\footnote{\url{https://www.tbench.ai/benchmarks/terminal-bench-2}} 把这个包含 89 个任务的子集划分为 4 个 easy、55 个 medium 和 30 个 hard 任务。

\paragraph{pass@1。}
对于在任务集 $D$ 上、每个任务做 $k$ 次 rollout 的某个配置，令 $r_{i,j} \in \{0, 1\}$ 表示第 $j$ 次 rollout 在任务 $i$ 上的二值奖励。pass@1 分数即为均值
\begin{equation}
    \mathrm{pass@1} = \frac{1}{k|D|} \sum_{i=1}^{|D|} \sum_{j=1}^{k} r_{i,j}.
\end{equation}
因基础设施异常（如沙箱崩溃或 API 超时）而终止的试验，不会被剔除，而是计为 $r=0$；这在严格意义上是比丢弃失败试验更苛刻的约定，能让我们的数字与官方 terminal-bench 排行榜保持可比。rollout 次数 $k$ 在不同实验中有所不同；每张表都会明确标注。

\paragraph{Token 成本与 Succ/Mtok。}
在 token 成本方面，我们把整个 rollout 中每一次 LLM 调用的 prompt 与 completion 相加计数，并以千为单位报告在完成试验上的均值，记为 Tokens k；因基础设施中止的试验被排除在外，以避免出现被截断的数字。为了比较那些以精度换成本的配置，我们通过下式把两者结合起来
\begin{equation}
    \mathrm{Succ/Mtok} = \frac{\mathrm{pass@1} \times 10^{6}}{\mathrm{mean\ tokens\ per\ trial}},
\end{equation}
即每百万 token 的期望成功次数。正文分别报告 pass@1 与 Tokens k，以保证每一个维度都清晰可读；表~\ref{tab:swe-verified-succ-mtok} 则把它们折算为 SWE-bench-verified 上按代码仓库划分的 Succ/Mtok，由表~\ref{tab:swe-verified-analysis} 的 pass@1 与 Tokens k 两列推导得到。

\begin{table}[t]
    \centering
    \small
    \setlength{\tabcolsep}{6pt}
    \renewcommand{\arraystretch}{1.05}
    \caption{SWE-bench-verified 上的成本效率，以 Succ/Mtok 报告，即每百万 token 的期望成功数。数值由表~\ref{tab:swe-verified-analysis} 按 $\mathrm{pass@1} \times 10^{3} / \text{Tokens k}$ 推导得到。越高越好。每行粗体标出最优。}
    \label{tab:swe-verified-succ-mtok}
    \begin{tabular}{lccccc}
        \toprule
        仓库 & $N$ & ACE & TF-GRPO & NexAU\textsubscript{0} & \textbf{AHE} \\
        \midrule
        全部            & 500 & 1.10 & 1.27 & 1.43 & \textbf{1.64} \\
        \midrule
        django         & 231 & 1.12 & 1.35 & 1.50 & \textbf{1.67} \\
        sympy          &  75 & 1.15 & 1.19 & 1.43 & \textbf{1.48} \\
        sphinx-doc     &  44 & 0.62 & 0.78 & 0.93 & \textbf{1.07} \\
        matplotlib     &  34 & 1.14 & 1.33 & 1.51 & \textbf{1.88} \\
        scikit-learn   &  32 & 2.08 & 2.48 & 3.06 & \textbf{3.40} \\
        pydata         &  22 & 1.37 & 1.50 & 2.00 & \textbf{2.15} \\
        astropy        &  22 & 1.08 & 1.26 & 0.82 & \textbf{1.81} \\
        \bottomrule
    \end{tabular}
\end{table}

